\documentclass[12pt,a4paper]{article}
\usepackage[UTF8]{ctex}
\usepackage{amsmath,amssymb,amsthm}
\usepackage{geometry}
\usepackage{graphicx}
\usepackage{hyperref}
\usepackage{booktabs}
\usepackage{array}
\usepackage{longtable}

\geometry{left=2.5cm,right=2.5cm,top=2.5cm,bottom=2.5cm}

\title{电机轴三维振动到像素平面位移的数学映射：证明过程}
\author{}
\date{}

\begin{document}

\maketitle

\tableofcontents

\newpage

\section{坐标系与变换}

\subsection{电机坐标系到相机坐标系的刚体变换}

设电机坐标系 $\{M\}$ 到相机坐标系 $\{C\}$ 的变换为旋转矩阵 $\mathbf{R}_{MC} \in \mathrm{SO}(3)$ 与平移向量 $\mathbf{t}_{MC} \in \mathbb{R}^3$：

\begin{equation}
\mathbf{P}^{(C)} = \mathbf{R}_{MC} \, \mathbf{P}^{(M)} + \mathbf{t}_{MC}
\label{eq:rigid_transform}
\end{equation}

其中 $\mathbf{R}_{MC}$ 满足 $\mathbf{R} \mathbf{R}^{\top} = \mathbf{I}$，$\det(\mathbf{R}) = 1$。

\subsection{旋转矩阵的Euler角参数化 (Z-Y-X)}

\begin{equation}
\mathbf{R}_{MC} = \mathbf{R}_Z(\gamma) \, \mathbf{R}_Y(\beta) \, \mathbf{R}_X(\alpha)
\label{eq:euler_rotation}
\end{equation}

展开为：

\begin{equation}
\mathbf{R}_{MC} =
\begin{bmatrix}
c_\gamma c_\beta & c_\gamma s_\beta s_\alpha - s_\gamma c_\alpha & c_\gamma s_\beta c_\alpha + s_\gamma s_\alpha \\
s_\gamma c_\beta & s_\gamma s_\beta s_\alpha + c_\gamma c_\alpha & s_\gamma s_\beta c_\alpha - c_\gamma s_\alpha \\
-s_\beta & c_\beta s_\alpha & c_\beta c_\alpha
\end{bmatrix}
\label{eq:rotation_matrix}
\end{equation}

其中 $c_\theta \equiv \cos\theta$，$s_\theta \equiv \sin\theta$。

\subsection{透视投影模型}

从相机坐标系 $\{C\}$ 到归一化像平面 $\{N\}$ 的映射：

\begin{equation}
\begin{bmatrix} x_n \\ y_n \end{bmatrix}
=
\frac{1}{Z_C}
\begin{bmatrix} X_C \\ Y_C \end{bmatrix}
\quad (Z_C > 0)
\label{eq:perspective_projection}
\end{equation}

\subsection{畸变模型}

\textbf{径向畸变}：

\begin{equation}
\begin{aligned}
x_{\text{rad}} &= x_n \, (1 + k_1 r^2 + k_2 r^4 + k_3 r^6) \\
y_{\text{rad}} &= y_n \, (1 + k_1 r^2 + k_2 r^4 + k_3 r^6)
\end{aligned}
\label{eq:radial_distortion}
\end{equation}

\textbf{切向畸变}：

\begin{equation}
\begin{aligned}
x_{\text{tan}} &= 2 p_1 x_n y_n + p_2 (r^2 + 2 x_n^2) \\
y_{\text{tan}} &= p_1 (r^2 + 2 y_n^2) + 2 p_2 x_n y_n
\end{aligned}
\label{eq:tangential_distortion}
\end{equation}

\textbf{畸变后坐标}：

\begin{equation}
\begin{aligned}
x_d &= x_{\text{rad}} + x_{\text{tan}} \\
y_d &= y_{\text{rad}} + y_{\text{tan}}
\end{aligned}
\label{eq:distorted_coordinates}
\end{equation}

其中 $r = \sqrt{x_n^2 + y_n^2}$。

\subsection{内参矩阵}

\begin{equation}
\begin{bmatrix} u \\ v \\ 1 \end{bmatrix}
=
\mathbf{K}
\begin{bmatrix} x_d \\ y_d \\ 1 \end{bmatrix}
\label{eq:intrinsic_matrix}
\end{equation}

其中：

\begin{equation}
\mathbf{K} =
\begin{bmatrix}
f_x & 0   & c_x \\
0   & f_y & c_y \\
0   & 0   & 1
\end{bmatrix}
\label{eq:K_matrix}
\end{equation}

\section{振动位移到像素位移的推导}

\subsection{振动位移定义}

电机轴上参考点 $\mathbf{P}_0 \in \mathbb{R}^3$ 在时刻 $t$ 的振动位移为 $\boldsymbol{\delta}(t)$：

\begin{equation}
\mathbf{P}^{(M)}(t) = \mathbf{P}_0^{(M)} + \boldsymbol{\delta}^{(M)}(t)
\label{eq:vibration_displacement}
\end{equation}

\subsection{振动在相机坐标系中的投影}

对公式~\eqref{eq:rigid_transform} 微分（$\mathbf{R}_{MC}$ 和 $\mathbf{t}_{MC}$ 为常数）：

\begin{equation}
\Delta \mathbf{P}^{(C)}(t) = \mathbf{R}_{MC} \, \boldsymbol{\delta}^{(M)}(t)
\label{eq:vibration_camera}
\end{equation}

展开为：

\begin{equation}
\begin{bmatrix}
\Delta X_C(t) \\ \Delta Y_C(t) \\ \Delta Z_C(t)
\end{bmatrix}
=
\begin{bmatrix}
R_{11} & R_{12} & R_{13} \\
R_{21} & R_{22} & R_{23} \\
R_{31} & R_{32} & R_{33}
\end{bmatrix}
\begin{bmatrix}
\delta_X(t) \\ \delta_Y(t) \\ \delta_Z(t)
\end{bmatrix}
\label{eq:vibration_camera_expanded}
\end{equation}

\subsection{透视投影映射的微分}

对 $x_n = X_C/Z_C$ 和 $y_n = Y_C/Z_C$ 进行全微分：

\begin{equation}
\begin{aligned}
\Delta x_n &= \frac{\partial}{\partial X_C}\!\left(\frac{X_C}{Z_C}\right) \Delta X_C
           + \frac{\partial}{\partial Z_C}\!\left(\frac{X_C}{Z_C}\right) \Delta Z_C \\
           &= \frac{1}{Z_C} \Delta X_C - \frac{X_C}{Z_C^2} \Delta Z_C \\[8pt]
\Delta y_n &= \frac{1}{Z_C} \Delta Y_C - \frac{Y_C}{Z_C^2} \Delta Z_C
\end{aligned}
\label{eq:projection_differential}
\end{equation}

用归一化坐标 $(x_n, y_n) = (X_C/Z_C, Y_C/Z_C)$ 简化：

\begin{equation}
\begin{aligned}
\Delta x_n &= \frac{1}{Z_C} \left( \Delta X_C - x_n \Delta Z_C \right) \\[6pt]
\Delta y_n &= \frac{1}{Z_C} \left( \Delta Y_C - y_n \Delta Z_C \right)
\end{aligned}
\label{eq:projection_differential_simplified}
\end{equation}

矩阵形式：

\begin{equation}
\begin{bmatrix} \Delta x_n \\ \Delta y_n \end{bmatrix}
=
\frac{1}{Z_C}
\begin{bmatrix}
1 & 0 & -x_n \\
0 & 1 & -y_n
\end{bmatrix}
\begin{bmatrix}
\Delta X_C \\ \Delta Y_C \\ \Delta Z_C
\end{bmatrix}
\label{eq:projection_matrix}
\end{equation}

\subsection{振动映射到归一化像平面位移}

将公式~\eqref{eq:vibration_camera_expanded} 代入公式~\eqref{eq:projection_matrix}：

\begin{equation}
\boxed{
\mathbf{\Delta} \mathbf{x}_n(t)
=
\frac{1}{Z_C}
\,
\mathbf{J}_{\Pi}(x_n, y_n)
\,
\mathbf{R}_{MC}
\,
\boldsymbol{\delta}^{(M)}(t)
}
\label{eq:vibration_to_normalized}
\end{equation}

其中 \textbf{投影 Jacobian} 矩阵 $\mathbf{J}_{\Pi} \in \mathbb{R}^{2 \times 3}$ 定义为：

\begin{equation}
\mathbf{J}_{\Pi}(x_n, y_n)
=
\begin{bmatrix}
1 & 0 & -x_n \\
0 & 1 & -y_n
\end{bmatrix}
\label{eq:projection_jacobian}
\end{equation}

\subsection{畸变对位移的影响}

对畸变映射 $(x_n, y_n) \mapsto (x_d, y_d)$ 进行全微分：

\begin{equation}
\begin{bmatrix} \Delta x_d \\ \Delta y_d \end{bmatrix}
=
\mathbf{J}_{\text{dist}}(x_n, y_n)
\begin{bmatrix} \Delta x_n \\ \Delta y_n \end{bmatrix}
\label{eq:distortion_differential}
\end{equation}

其中 $\mathbf{J}_{\text{dist}} \in \mathbb{R}^{2 \times 2}$ 为畸变模型的 Jacobian 矩阵：

\begin{equation}
\mathbf{J}_{\text{dist}} =
\begin{bmatrix}
\partial x_d / \partial x_n & \partial x_d / \partial y_n \\
\partial y_d / \partial x_n & \partial y_d / \partial y_n
\end{bmatrix}
\label{eq:distortion_jacobian}
\end{equation}

令 $r^2 = x_n^2 + y_n^2$，$\alpha_r = 1 + k_1 r^2 + k_2 r^4 + k_3 r^6$，$\alpha_r' = 2k_1 + 4k_2 r^2 + 6k_3 r^4$。

径向畸变 Jacobian：

\begin{equation}
\begin{aligned}
\frac{\partial (x_n \alpha_r)}{\partial x_n} &= \alpha_r + x_n^2 \alpha_r' \\
\frac{\partial (x_n \alpha_r)}{\partial y_n} &= x_n y_n \alpha_r' \\
\frac{\partial (y_n \alpha_r)}{\partial x_n} &= x_n y_n \alpha_r' \\
\frac{\partial (y_n \alpha_r)}{\partial y_n} &= \alpha_r + y_n^2 \alpha_r'
\end{aligned}
\label{eq:radial_jacobian}
\end{equation}

切向畸变 Jacobian：

\begin{equation}
\begin{aligned}
\frac{\partial x_{\text{tan}}}{\partial x_n} &= 2p_1 y_n + 4p_2 x_n & \frac{\partial x_{\text{tan}}}{\partial y_n} &= 2p_1 x_n + 2p_2 y_n \\
\frac{\partial y_{\text{tan}}}{\partial x_n} &= 2p_1 x_n + 2p_2 y_n & \frac{\partial y_{\text{tan}}}{\partial y_n} &= 4p_1 y_n + 2p_2 x_n
\end{aligned}
\label{eq:tangential_jacobian}
\end{equation}

则：

\begin{equation}
\mathbf{J}_{\text{dist}} =
\begin{bmatrix}
\alpha_r + x_n^2 \alpha_r' + 2p_1 y_n + 4p_2 x_n & x_n y_n \alpha_r' + 2p_1 x_n + 2p_2 y_n \\[4pt]
x_n y_n \alpha_r' + 2p_1 x_n + 2p_2 y_n & \alpha_r + y_n^2 \alpha_r' + 4p_1 y_n + 2p_2 x_n
\end{bmatrix}
\label{eq:distortion_jacobian_full}
\end{equation}

\subsection{从归一化畸变坐标到像素坐标}

对 $u = f_x x_d + c_x$ 和 $v = f_y y_d + c_y$ 微分（主点 $(c_x, c_y)$ 为常数）：

\begin{equation}
\begin{aligned}
\Delta u &= f_x \, \Delta x_d \\
\Delta v &= f_y \, \Delta y_d
\end{aligned}
\label{eq:pixel_differential}
\end{equation}

矩阵形式：

\begin{equation}
\begin{bmatrix} \Delta u \\ \Delta v \end{bmatrix}
=
\begin{bmatrix}
f_x & 0 \\
0   & f_y
\end{bmatrix}
\begin{bmatrix} \Delta x_d \\ \Delta y_d \end{bmatrix}
\label{eq:pixel_matrix}
\end{equation}

\subsection{完整映射的合成}

联立公式~\eqref{eq:vibration_to_normalized}、\eqref{eq:distortion_differential}、\eqref{eq:pixel_matrix}：

\begin{equation}
\boxed{
\begin{bmatrix} \Delta u(t) \\[2pt] \Delta v(t) \end{bmatrix}
=
\frac{1}{Z_C}
\,
\begin{bmatrix}
f_x & 0 \\
0   & f_y
\end{bmatrix}
\,
\mathbf{J}_{\text{dist}}(x_n, y_n)
\,
\mathbf{J}_{\Pi}(x_n, y_n)
\,
\mathbf{R}_{MC}
\,
\begin{bmatrix}
\delta_X(t) \\[2pt] \delta_Y(t) \\[2pt] \delta_Z(t)
\end{bmatrix}
}
\label{eq:complete_mapping}
\end{equation}

\subsection{灵敏度矩阵}

定义 \textbf{振动灵敏度矩阵} $\mathbf{S} \in \mathbb{R}^{2 \times 3}$：

\begin{equation}
\boxed{
\mathbf{S}(\mathbf{P}_0^{(M)}; \mathbf{R}_{MC}, \mathbf{t}_{MC}, \mathbf{K}, \mathbf{d})
=
\frac{1}{Z_C}
\,
\mathbf{K}_{2\times 2}
\,
\mathbf{J}_{\text{dist}}(x_n, y_n)
\,
\mathbf{J}_{\Pi}(x_n, y_n)
\,
\mathbf{R}_{MC}
}
\label{eq:sensitivity_matrix}
\end{equation}

使得：

\begin{equation}
\boxed{
\mathbf{\Delta u}(t) = \mathbf{S} \cdot \boldsymbol{\delta}^{(M)}(t)
}
\label{eq:sensitivity_relation}
\end{equation}

其中 $\mathbf{d} = (k_1, k_2, k_3, p_1, p_2)$ 为畸变参数向量。

\section{椭圆投影与振动参数反演}

\subsection{电机轴圆截面的数学描述}

在电机坐标系 $\{M\}$ 中，$Z_M = z_0$ 处的横截面为：

\begin{equation}
\mathcal{C}^{(M)} = \left\{
\mathbf{P}^{(M)} \in \mathbb{R}^3 \,\Big|\,
X_M^2 + Y_M^2 = r^2,\; Z_M = z_0
\right\}
\label{eq:circle_cross_section}
\end{equation}

参数化形式：

\begin{equation}
\mathbf{P}^{(M)}(\theta) =
\begin{bmatrix}
r \cos\theta \\[2pt]
r \sin\theta \\[2pt]
z_0
\end{bmatrix},
\quad \theta \in [0, 2\pi)
\label{eq:circle_parameterization}
\end{equation}

\subsection{圆截面在相机坐标系中的表示}

利用公式~\eqref{eq:rigid_transform} 的刚体变换：

\begin{equation}
\mathbf{P}^{(C)}(\theta) =
\mathbf{R}_{MC} \, \mathbf{P}^{(M)}(\theta) + \mathbf{t}_{MC}
\label{eq:circle_camera}
\end{equation}

展开为：

\begin{equation}
\begin{aligned}
X_C(\theta) &= R_{11} r \cos\theta + R_{12} r \sin\theta + R_{13} z_0 + t_X \\
Y_C(\theta) &= R_{21} r \cos\theta + R_{22} r \sin\theta + R_{23} z_0 + t_Y \\
Z_C(\theta) &= R_{31} r \cos\theta + R_{32} r \sin\theta + R_{33} z_0 + t_Z
\end{aligned}
\label{eq:circle_camera_expanded}
\end{equation}

定义圆心在相机坐标系中的坐标：

\begin{equation}
\begin{bmatrix}
X_{C0} \\ Y_{C0} \\ Z_{C0}
\end{bmatrix}
=
\begin{bmatrix}
R_{13} z_0 + t_X \\
R_{23} z_0 + t_Y \\
R_{33} z_0 + t_Z
\end{bmatrix}
\label{eq:circle_center_camera}
\end{equation}

则公式~\eqref{eq:circle_camera_expanded} 可写为：

\begin{equation}
\begin{bmatrix}
X_C(\theta) \\ Y_C(\theta) \\ Z_C(\theta)
\end{bmatrix}
=
\begin{bmatrix}
X_{C0} \\ Y_{C0} \\ Z_{C0}
\end{bmatrix}
+
r
\begin{bmatrix}
R_{11} \cos\theta + R_{12} \sin\theta \\
R_{21} \cos\theta + R_{22} \sin\theta \\
R_{31} \cos\theta + R_{32} \sin\theta
\end{bmatrix}
\label{eq:circle_camera_simplified}
\end{equation}

\subsection{椭圆投影方程的推导}

由透视投影公式~\eqref{eq:perspective_projection}，归一化像平面坐标为：

\begin{equation}
\begin{bmatrix}
x_n(\theta) \\[2pt] y_n(\theta)
\end{bmatrix}
=
\frac{1}{Z_C(\theta)}
\begin{bmatrix}
X_C(\theta) \\[2pt] Y_C(\theta)
\end{bmatrix}
\label{eq:ellipse_projection}
\end{equation}

根据 $2r \ll Z_C$，对 $Z_C(\theta)$ 做一阶展开：

\begin{equation}
Z_C(\theta) = Z_{C0} + r \,(R_{31} \cos\theta + R_{32} \sin\theta) + \mathcal{O}(r^2/Z_{C0})
\label{eq:depth_expansion}
\end{equation}

对公式~\eqref{eq:ellipse_projection} 做一阶 Taylor 展开：

\begin{equation}
\begin{aligned}
x_n(\theta) &\approx \frac{X_{C0}}{Z_{C0}} +
\frac{r}{Z_{C0}}(R_{11} \cos\theta + R_{12} \sin\theta)
- \frac{X_{C0}}{Z_{C0}}\cdot\frac{r}{Z_{C0}}(R_{31} \cos\theta + R_{32} \sin\theta) \\[6pt]
y_n(\theta) &\approx \frac{Y_{C0}}{Z_{C0}} +
\frac{r}{Z_{C0}}(R_{21} \cos\theta + R_{22} \sin\theta)
- \frac{Y_{C0}}{Z_{C0}}\cdot\frac{r}{Z_{C0}}(R_{31} \cos\theta + R_{32} \sin\theta)
\end{aligned}
\label{eq:ellipse_approximation}
\end{equation}

定义圆心在归一化像平面上的投影坐标 $(x_{n0}, y_{n0})$：

\begin{equation}
x_{n0} = \frac{X_{C0}}{Z_{C0}}, \quad
y_{n0} = \frac{Y_{C0}}{Z_{C0}}
\label{eq:circle_center_normalized}
\end{equation}

以及缩放系数 $\alpha = r / Z_{C0}$（无量纲，$0 < \alpha \ll 1$）。则公式~\eqref{eq:ellipse_approximation} 化为：

\begin{equation}
\begin{aligned}
x_n(\theta) &\approx x_{n0} +
\alpha\big[\,(R_{11} - x_{n0} R_{31})\cos\theta + (R_{12} - x_{n0} R_{32})\sin\theta\,\big] \\[6pt]
y_n(\theta) &\approx y_{n0} +
\alpha\big[\,(R_{21} - y_{n0} R_{31})\cos\theta + (R_{22} - y_{n0} R_{32})\sin\theta\,\big]
\end{aligned}
\label{eq:ellipse_final}
\end{equation}

\subsection{椭圆参数化标准型}

将公式~\eqref{eq:ellipse_final} 写成向量形式。定义向量 $\mathbf{p} = (x_n, y_n)^{\top}$，$\mathbf{p}_0 = (x_{n0}, y_{n0})^{\top}$，则：

\begin{equation}
\mathbf{p}(\theta) = \mathbf{p}_0 + \alpha \,
\mathbf{A}
\begin{bmatrix}
\cos\theta \\ \sin\theta
\end{bmatrix}
\label{eq:ellipse_vector}
\end{equation}

其中 $\mathbf{A} \in \mathbb{R}^{2 \times 2}$ 为\textbf{形变矩阵}：

\begin{equation}
\boxed{
\mathbf{A} =
\begin{bmatrix}
R_{11} - x_{n0} R_{31} & R_{12} - x_{n0} R_{32} \\[4pt]
R_{21} - y_{n0} R_{31} & R_{22} - y_{n0} R_{32}
\end{bmatrix}
}
\label{eq:deformation_matrix}
\end{equation}

\textbf{证明}：公式~\eqref{eq:ellipse_vector} 中 $(\cos\theta, \sin\theta)^{\top}$ 参数化了一个半径为 1 的单位圆。乘以矩阵 $\mathbf{M} = \alpha \mathbf{A}$ 后，单位圆变为一个中心在原点的椭圆，其形状由 $\mathbf{M}$ 的奇异值决定。最后平移 $\mathbf{p}_0$ 得到中心在 $\mathbf{p}_0$ 的椭圆，证毕。

\subsection{椭圆几何参数与电机轴-成像平面夹角的关系}

椭圆由形变矩阵 $\mathbf{M} = \alpha \mathbf{A}$ 的\textbf{奇异值分解 (SVD)} 完全刻画：

\begin{equation}
\mathbf{M} = \mathbf{U} \boldsymbol{\Sigma} \mathbf{V}^{\top}
\label{eq:svd}
\end{equation}

其中 $\boldsymbol{\Sigma} = \operatorname{diag}(\sigma_1, \sigma_2)$，$\sigma_1 \ge \sigma_2 \ge 0$。

\subsection{仅考虑一个旋转角的简化情况}

设电机坐标系 $\{M\}$ 到相机坐标系 $\{C\}$ 的变换仅包含绕 $X_M$ 轴旋转 $\alpha$ 角：

\begin{equation}
\mathbf{R}_{MC} =
\begin{bmatrix}
1 & 0 & 0 \\
0 & \cos\alpha & -\sin\alpha \\
0 & \sin\alpha & \cos\alpha
\end{bmatrix}
\label{eq:simple_rotation}
\end{equation}

假设圆的投影中心近似位于主点附近（$x_{n0} \approx 0, y_{n0} \approx 0$），则形变矩阵 $\mathbf{A}$ 简化为：

\begin{equation}
\mathbf{A} \approx
\begin{bmatrix}
1 & 0 \\
0 & \cos\alpha
\end{bmatrix}
\label{eq:simple_deformation}
\end{equation}

代入公式~\eqref{eq:ellipse_vector}：

\begin{equation}
\begin{aligned}
x_n(\theta) &\approx \alpha \cos\theta \\[4pt]
y_n(\theta) &\approx \alpha \cos\alpha \sin\theta
\end{aligned}
\;\Longrightarrow\;
\frac{x_n^2}{\alpha^2} + \frac{y_n^2}{\alpha^2 \cos^2\alpha} = 1
\label{eq:ellipse_equation}
\end{equation}

\textbf{电机轴与成像平面的夹角} $\psi$ 与椭圆偏心率的\textbf{显式关系}：

\begin{equation}
\boxed{
|\sin\psi| = e = \sqrt{1 - \frac{b^2}{a^2}}
\quad\Longrightarrow\quad
\psi = \arcsin\!\left(\sqrt{1 - \frac{b^2}{a^2}}\right)
}
\label{eq:angle_eccentricity}
\end{equation}

其中 $\psi \equiv \pi/2 - \alpha$ 为电机轴 $Z_M$ 与成像平面（即 $X_CY_C$ 平面）的法线 $Z_C$ 之间的夹角。

\section{振动参数的三维反演}

\subsection{时变椭圆模型}

当电机轴发生振动时，圆截面在电机坐标系中的位置随时间变化：

\begin{equation}
\mathbf{P}^{(M)}(\theta, t) =
\begin{bmatrix}
r\cos\theta + \delta_X(t) \\[2pt]
r\sin\theta + \delta_Y(t) \\[2pt]
z_0 + \delta_Z(t)
\end{bmatrix}
\label{eq:time_varying_circle}
\end{equation}

\subsection{椭圆中心位移}

椭圆中心在归一化像平面上的时变投影为：

\begin{equation}
\mathbf{p}_0(t) = \mathbf{p}_0(0) + \frac{1}{Z_{C0}} \,
\mathbf{J}_{\Pi}(x_{n0}, y_{n0}) \,
\mathbf{R}_{MC} \,
\boldsymbol{\delta}^{(M)}(t)
\label{eq:ellipse_center_motion}
\end{equation}

在像素坐标系中（忽略畸变），椭圆中心的像素位移为：

\begin{equation}
\boxed{
\begin{bmatrix}
\Delta u_c(t) \\[2pt] \Delta v_c(t)
\end{bmatrix}
=
\frac{1}{Z_{C0}}
\begin{bmatrix}
f_x & 0 \\
0   & f_y
\end{bmatrix}
\begin{bmatrix}
1 & 0 & -x_{n0} \\
0 & 1 & -y_{n0}
\end{bmatrix}
\mathbf{R}_{MC}
\begin{bmatrix}
\delta_X(t) \\[2pt] \delta_Y(t) \\[2pt] \delta_Z(t)
\end{bmatrix}
}
\label{eq:ellipse_center_pixel}
\end{equation}

\subsection{椭圆形变}

椭圆形状（长短轴）随振动位移的变化来源于深度 $Z_{C0}$ 的变化。设深度为 $Z_{C0}(t) = Z_{C0}^{(0)} + \Delta Z_C(t)$，则时变半长轴和半短轴为：

\begin{equation}
a(t) = \frac{r}{Z_{C0}(t)} \cdot \sigma_1(\mathbf{A})
\approx a_0 \left(1 - \frac{\Delta Z_C(t)}{Z_{C0}^{(0)}}\right)
\label{eq:time_varying_semi_major}
\end{equation}

\begin{equation}
b(t) = \frac{r}{Z_{C0}(t)} \cdot \sigma_2(\mathbf{A})
\approx b_0 \left(1 - \frac{\Delta Z_C(t)}{Z_{C0}^{(0)}}\right)
\label{eq:time_varying_semi_minor}
\end{equation}

其中 $\Delta Z_C(t) = R_{31}\delta_X(t) + R_{32}\delta_Y(t) + R_{33}\delta_Z(t)$。

\subsection{观测模型}

从每一帧图像中提取椭圆参数，构成观测量向量：

\begin{equation}
\mathbf{y}(t) =
\begin{bmatrix}
\Delta u_c(t) \\[2pt] \Delta v_c(t) \\[2pt]
\Delta a(t) \\[2pt] \Delta b(t) \\[2pt]
\Delta \phi(t)
\end{bmatrix}
\in \mathbb{R}^5
\label{eq:observation_vector}
\end{equation}

建立从振动状态向量 $\boldsymbol{\delta}^{(M)}(t)$ 到观测量 $\mathbf{y}(t)$ 的线性化映射：

\begin{equation}
\mathbf{y}(t) = \mathbf{H} \cdot \boldsymbol{\delta}^{(M)}(t) + \mathbf{n}(t)
\label{eq:observation_model}
\end{equation}

其中 $\mathbf{H} \in \mathbb{R}^{5 \times 3}$ 为\textbf{观测矩阵}。

\subsection{三轴振动分离的封闭形式}

在简化配置（投影中心靠近主点）下：

\begin{equation}
\Delta u_c(t) \approx \frac{f_x}{Z_{C0}} \,\delta_X(t)
\label{eq:u_simplified}
\end{equation}

\begin{equation}
\Delta v_c(t) \approx \frac{f_y}{Z_{C0}} \left(\cos\alpha \, \delta_Y(t) - \sin\alpha \, \delta_Z(t)\right)
\label{eq:v_simplified}
\end{equation}

\begin{equation}
\frac{\Delta a(t)}{a_0} \approx -\frac{1}{Z_{C0}} \left(\sin\alpha \, \delta_Y(t) + \cos\alpha \, \delta_Z(t)\right)
\label{eq:a_simplified}
\end{equation}

联立公式~\eqref{eq:u_simplified}-\eqref{eq:a_simplified} 构成三个方程三个未知数 $(\delta_X, \delta_Y, \delta_Z)$，解得：

\begin{equation}
\boxed{
\delta_X(t) = \frac{Z_{C0}}{f_x} \, \Delta u_c(t)
}
\label{eq:delta_X}
\end{equation}

\begin{equation}
\boxed{
\delta_Y(t) = \frac{Z_{C0}}{f_y} \cos\alpha \, \Delta v_c(t)
- Z_{C0} \sin\alpha \, \frac{\Delta a(t)}{a_0}
}
\label{eq:delta_Y}
\end{equation}

\begin{equation}
\boxed{
\delta_Z(t) = -\frac{Z_{C0}}{f_y} \sin\alpha \, \Delta v_c(t)
- Z_{C0} \cos\alpha \, \frac{\Delta a(t)}{a_0}
}
\label{eq:delta_Z}
\end{equation}

\textbf{推导说明}：将公式~\eqref{eq:v_simplified}-\eqref{eq:a_simplified} 视为关于 $(\delta_Y, \delta_Z)$ 的 2×2 线性方程组：

\begin{equation}
\begin{bmatrix}
\cos\alpha & -\sin\alpha \\
\sin\alpha & \cos\alpha
\end{bmatrix}
\begin{bmatrix} \delta_Y \\ \delta_Z \end{bmatrix}
=
\begin{bmatrix}
(Z_{C0}/f_y) \Delta v_c \\
-Z_{C0} (\Delta a / a_0)
\end{bmatrix}
\label{eq:linear_system}
\end{equation}

系数矩阵为旋转矩阵，行列式 $= \cos^2\alpha + \sin^2\alpha = 1$，其逆矩阵即为转置。由此解得上述公式。

\subsection{频率信息提取}

对公式~\eqref{eq:delta_X}-\eqref{eq:delta_Z} 估计的时序 $\hat{\delta}_X(t_k), \hat{\delta}_Y(t_k), \hat{\delta}_Z(t_k)$ 进行傅里叶变换：

\begin{equation}
\tilde{\delta}_i(f) = \sum_{k=0}^{N-1} \hat{\delta}_i(t_k) e^{-j2\pi f k / N}
\label{eq:fourier_transform}
\end{equation}

各频率分量的幅度和初相由峰值提取：

\begin{equation}
A_{i,k} = 2|\tilde{\delta}_i(f_k)|, \quad
\varphi_{i,k} = \arg(\tilde{\delta}_i(f_k))
\label{eq:frequency_extraction}
\end{equation}

\section{误差传播分析}

设椭圆参数提取误差的协方差矩阵为 $\boldsymbol{\Sigma}_{\mathbf{y}} \in \mathbb{R}^{5 \times 5}$，则振动估计误差的协方差矩阵为：

\begin{equation}
\boldsymbol{\Sigma}_{\boldsymbol{\delta}} = (\mathbf{H}^{\top} \mathbf{H})^{-1} \mathbf{H}^{\top} \boldsymbol{\Sigma}_{\mathbf{y}} \mathbf{H} (\mathbf{H}^{\top} \mathbf{H})^{-1}
\label{eq:error_covariance}
\end{equation}

各反演参数的\textbf{不确定度}为：

\begin{equation}
\sigma_{\delta_i} = \sqrt{[\boldsymbol{\Sigma}_{\boldsymbol{\delta}}]_{ii}}, \quad i \in \{X, Y, Z\}
\label{eq:uncertainty}
\end{equation}

\section{变量定义}

\begin{longtable}{p{4cm}p{6cm}p{3cm}}
\toprule
\textbf{符号} & \textbf{物理含义} & \textbf{单位} \\
\midrule
$\boldsymbol{\delta}^{(M)} = (\delta_X, \delta_Y, \delta_Z)^{\top}$ & 电机轴三轴振动位移 & m \\
$\mathbf{P}_0^{(M)} = (X_0, Y_0, Z_0)^{\top}$ & 观测点在电机坐标系中的平衡位置 & m \\
$\mathbf{R}_{MC}$ & 电机系→相机系旋转矩阵 & 无量纲 \\
$\mathbf{t}_{MC} = (t_X, t_Y, t_Z)^{\top}$ & 电机系→相机系平移向量 & m \\
$Z_C$ & 观测点在相机系中的深度 & m \\
$f_x, f_y$ & 等效焦距（像素单位） & pixel \\
$c_x, c_y$ & 主点像素坐标 & pixel \\
$k_1, k_2, k_3$ & 径向畸变系数 & 无量纲 \\
$p_1, p_2$ & 切向畸变系数 & 无量纲 \\
$(x_n, y_n)$ & 像点的归一化坐标 & 无量纲 \\
$r = \sqrt{x_n^2 + y_n^2}$ & 归一化径向距离 & 无量纲 \\
$\mathbf{J}_{\Pi}$ & 投影 Jacobian ($2 \times 3$) & 无量纲 \\
$\mathbf{J}_{\text{dist}}$ & 畸变 Jacobian ($2 \times 2$) & 无量纲 \\
$\mathbf{S}$ & 振动-像素灵敏度矩阵 ($2 \times 3$) & pixel/m \\
$(\Delta u, \Delta v)^{\top}$ & 像素位移 & pixel \\
$r$ & 电机轴横截面半径 & m \\
$z_0$ & 观测截面在电机轴上的轴向位置 & m \\
$\mathbf{A}$ & 形变矩阵 ($2 \times 2$) & 无量纲 \\
$\alpha = r / Z_{C0}$ & 缩放系数 & 无量纲 \\
$a, b$ & 椭圆半长轴、半短轴（归一化坐标） & 无量纲 \\
$e$ & 椭圆偏心率 & 无量纲 \\
$\psi$ & 电机轴与成像平面法线夹角 & rad \\
$\phi$ & 椭圆长轴方向角 & rad \\
$\mathbf{y}(t)$ & 观测量向量 ($5 \times 1$) & 混合 \\
$\mathbf{H}$ & 观测矩阵 ($5 \times 3$) & 混合 \\
$\sigma_1, \sigma_2$ & 形变矩阵 $\mathbf{M}$ 的奇异值 & 无量纲 \\
$A_{i,k}, f_{i,k}, \varphi_{i,k}$ & 振动分量的振幅、频率、初相 & m, Hz, rad \\
\bottomrule
\end{longtable}

\section{核心结论}

\subsection{结论一：振动-像素映射的线性化表达式}

电机轴三维振动 $\boldsymbol{\delta}^{(M)}(t)$ 到像素位移 $\mathbf{\Delta u}(t)$ 的映射为：

\begin{equation}
\boxed{\mathbf{\Delta u}(t) = \mathbf{S} \cdot \boldsymbol{\delta}^{(M)}(t)}
\label{eq:conclusion_1}
\end{equation}

其中灵敏度矩阵 $\mathbf{S} = \frac{1}{Z_C} \mathbf{K}_{2\times 2} \mathbf{J}_{\text{dist}} \mathbf{J}_{\Pi} \mathbf{R}_{MC}$，表明像素位移与振动位移成\textbf{线性关系}。

\subsection{结论二：椭圆偏心率反映轴-平面夹角}

电机轴圆截面在透视投影下成像为椭圆，椭圆偏心率 $e$ 与电机轴-成像平面夹角 $\psi$ 满足：

\begin{equation}
\boxed{\psi = \arcsin\!\left(\sqrt{1 - \frac{b^2}{a^2}}\right)}
\label{eq:conclusion_2}
\end{equation}

\begin{itemize}
\item $e = 0$（圆）$\Rightarrow$ 轴与平面平行
\item $e \to 1$（扁椭圆）$\Rightarrow$ 轴趋于垂直
\end{itemize}

\subsection{结论三：单目相机可分离三轴振动}

利用椭圆中心位移 $(\Delta u_c, \Delta v_c)$ 和椭圆尺度变化 $\Delta a/a_0$，可\textbf{封闭求解}三轴振动：

\begin{equation}
\boxed{
\begin{aligned}
\delta_X &= \frac{Z_{C0}}{f_x} \Delta u_c \\[6pt]
\delta_Y &= \frac{Z_{C0}}{f_y} \cos\alpha \cdot \Delta v_c - Z_{C0} \sin\alpha \cdot \frac{\Delta a}{a_0} \\[6pt]
\delta_Z &= -\frac{Z_{C0}}{f_y} \sin\alpha \cdot \Delta v_c - Z_{C0} \cos\alpha \cdot \frac{\Delta a}{a_0}
\end{aligned}
}
\label{eq:conclusion_3}
\end{equation}

\textbf{物理意义}：单目相机通过椭圆的\textbf{位置变化}（2个观测量）和\textbf{形状变化}（1个观测量），共3个独立观测量，恰好可解3个未知振动分量。

\subsection{总结}

\begin{center}
\fbox{\parbox{0.85\textwidth}{
\centering\Large\textbf{单目相机 + 圆截面先验知识 $\Rightarrow$ 完整的三维振动测量}

\vspace{0.5em}

\normalsize 不需要多相机、激光测振仪或加速度传感器，仅通过单目视觉对电机轴截面椭圆的观测，即可实现三轴振动的实时反演。
}}
\end{center}

\end{document}